\documentclass{article}
\usepackage[UTF8]{ctex}
\usepackage{amsmath,amssymb,amsfonts}
\usepackage{enumitem}
\usepackage{geometry}
\geometry{a4paper, margin=1in}
% === 错题本排版优化（不参与题目标识） ===
\setlength{\parindent}{0pt}
\setlength{\parskip}{0.35em plus 0.10em minus 0.08em}
\setlength{\emergencystretch}{3em}
\setlist{itemsep=0.25em, topsep=0.40em, parsep=0pt}
\allowdisplaybreaks[2]
% === 错题本排版优化结束 ===
\begin{document}
\section*{中国人民大学2026年数学分析真题}
\section*{21．中国人民大学2026年数学分析试卷}

\subsection{题目1}
1．（15 分）讨论 \(f(x)=\sin \sqrt{x}\) 和 \(g(x)=\sin \left(x^{2}\right)\) 在 \([0,+\infty)\) 上是否一致连续．

\subsection{题目2}
2．（15 分）设 \(a_{n}>0, n=1,2, \cdots\) ，且 \(\varlimsup_{n \rightarrow \infty} \frac{a_{n+1}}{a_{n}}<+\infty\) ，证明：\(\varlimsup_{n \rightarrow \infty} \sqrt[n]{a_{n}} \leq \varlimsup_{n \rightarrow \infty} \frac{a_{n+1}}{a_{n}}\) ．

\subsection{题目3}
3．（15 分）证明：对任意的 \(x \in[0,1]\) ；有 \(1+x^{2} \leq 2^{x}\) ．

\subsection{题目4}
4．（15 分）设 \(f(x)\) 为 \(\mathbb{R}\) 上周期为 \(\pi\) 的函数，且在 \(\left[-\frac{\pi}{2}, \frac{\pi}{2}\right)\) 上，有

\[
f(x)= \begin{cases}\frac{\pi}{2}+x, & x \in\left[-\frac{\pi}{2}, 0\right) \\ \frac{\pi}{2}-x, & x \in\left[0, \frac{\pi}{2}\right)\end{cases}
\]

请将 \(f(x)\) 在 \(\left[-\frac{\pi}{2}, \frac{\pi}{2}\right)\) 上展开为傅里叶级数，并由此求 \(\sum_{n=1}^{\infty} \frac{1}{(2 n-1)^{2}}\) ．

\subsection{题目5}
5．（15 分）设对于任意的正整数 \(n, f_{n}(x)\) 均为定义在 \([a, b]\) 上的黎曼可积函数，且 \(\left\{f_{n}(x)\right\}\) 在 \([a, b]\) 上一致收敛于函数 \(f(x)\) 。证明：\(f(x)\) 是 \([a, b]\) 上的黎曼可积函数。

\subsection{题目6}
6．（15 分）计算积分 \(\int_{0}^{\frac{\pi}{2}} \frac{1}{\sin \theta} \ln \frac{1+t \sin \theta}{1-t \sin \theta} \mathrm{~d} \theta,|t|<1\) ．

\subsection{题目7}
7．（15 分）设

\[
f(x, y)= \begin{cases}\frac{x^{2} y+x^{4}}{x^{2}+y^{2}}, & x^{2}+y^{2} \neq 0 \\ 0, & x^{2}+y^{2}=0\end{cases}
\]

讨论 \(f(x, y)\) 在 \((0,0)\) 的连续性和可微性．又设 \(x(t)=t, y(t)=t^{2}\) ，讨论复合函数 \(f(x(t), y(t))\) 在\(t=0\) 的导数是否满足链式法则．

\subsection{题目8}
8．（15 分）设 \(f(x, y)\) 为定义在 \(D=\left\{(x, y) \mid x^{2}+y^{2}<1\right\}\) 上的二元函数，在 \(D\) 内对 \(x\) 和 \(y\) 具有直到二阶的连续偏导数，且 \(\left(x_{0}, y_{0}\right) \in D\) 是 \(f(x, y)\) 的极小值点。证明：

\[
f_{x x}\left(x_{0}, y_{0}\right)+2 f_{x y}\left(x_{0}, y_{0}\right)+f_{y y}\left(x_{0}, y_{0}\right) \geq 0
\]

\subsection{题目9}
9．（15 分）设反常积分

\[
I=\iiint_{D} \frac{x^{2}}{\left(x^{2}+y^{2}+z^{2}\right)^{2}} \mathrm{~d} x \mathrm{~d} y \mathrm{~d} z .
\]

其中 \(D=\left\{(x, y, z) \mid x^{2}+y^{2}+z^{2} \leq 4 x\right\}\) 。证明反常积分 \(I\) 收敛，并求出 \(I\) 的值。

\subsection{题目10}
10．（15 分）设 \(D\) 是平面区域，\(u(x, y)\) 为定义在 \(D\) 上的二元函数，且在 \(D\) 内对 \(x\) 和 \(y\) 具有直到二阶的连续偏导数。证明：\(u(x, y)\) 在 \(D\) 上满足 \(\frac{\partial^{2} u}{\partial x^{2}}+\frac{\partial^{2} u}{\partial y^{2}}=0\) 的充要条件为对于任意 \(D\) 内的圆周 \(L\) ，且\(L\) 所围圆 \(O\) 含于 \(D\) ，有 \(\oint_{L} \frac{\partial u}{\partial \mathbf{n}} \mathrm{~d} s=0\) ，其中 \(\mathbf{n}\) 取圆周 \(L\) 的外法向。
\end{document}
